\section*{Chapter \ref{ch:m1:options-markets-europe-asia} --- Options Markets in Europe and Asia}

\begin{solution}{exo:m1:options-markets-europe-asia:1}
Notional $25 \times 22\,000 \times 0.012 = \$6\,600$; premium $25 \times 35 \times
0.012 = \$10.50$.
\end{solution}

\begin{solution}{exo:m1:options-markets-europe-asia:2}
Notional \$560\,000; premium \$1\,600. One such contract carries the premium
of 152 of the small ones, and the notional of 85.
\end{solution}

\begin{solution}{exo:m1:options-markets-europe-asia:3}
30 billion. With 24 billion observed, notional activity is $24 \times 75/(90
\times 25) = 80\%$ of what it was: a fifth of the activity left with the
smaller participants.
\end{solution}

\begin{solution}{exo:m1:options-markets-europe-asia:4}
2024: $170.2 + 19.8 + 10.0 + 4.2 + 2.3 = 206.4$ billion; Asia-Pacific 82\%.
2025: $75.59/119.29 = 63\%$. Check: $1 - 119.29/206.4 = 42.2\%$.
\end{solution}

\begin{solution}{exo:m1:options-markets-europe-asia:5}
M1: $945 \times 4\% \times \tfrac14 = \$9.45$ billion. M2: $1\,360 \times 1.5\%
\times \tfrac14 = \$5.1$ billion. M1's pool is larger despite its smaller
\omterm{def:m1:options-markets-europe-asia:turnover}{premium turnover}, because spreads are wide. Before entering: fees and taxes;
who may trade and in what legal form; how much of the pool incumbents with
better access already take; the \omterm{def:m1:what-a-trading-firm-does:adverse-selection}{adverse selection} in each flow; and how
long the contract specification will stay as it is.
\end{solution}

\begin{solution}{exo:m1:options-markets-europe-asia:6}
The mover must sustain the pressure for thirty minutes instead of one print,
against arbitrageurs who have time to lean the other way; and each minute's
price counts for a thirtieth. Cost: nobody can trade \emph{at} the
settlement value, so hedgers who want to exit at settlement must execute a
schedule over the window and bear \omterm{def:m1:the-buy-side:benchmark}{tracking error} against the average.
\end{solution}

\begin{solution}{exo:m1:options-markets-europe-asia:7}
Before: contracts 98.1 / 0.9 / 0.4 / 0.5\%; notional 48.3 / 38.7 / 1.6 /
11.4\%; premium 38.2 / 54.9 / 6.9 / 0\%. After (M1 at 24 billion contracts of
lot 75): M1's share of contracts falls to about 93\%, still overwhelming,
while its notional and premium, each 80\% of before, give shares of about 43\%
and 33\%. The contract table hardly notices a change that removed a fifth of
the market's activity.
\end{solution}

\begin{solution}{exo:m1:options-markets-europe-asia:8}
The combination. A regulator, and a court, will look at the firm as one
entity: a large short index-options position whose payoff depends on the
closing index, held while the same firm's trades in the constituents move
that index, first up while the options are being sold, then down into the
settlement. Each leg's local rationale does not answer the question of
intent for the whole, and the firm's own records show the whole. What is
missing is a control that aggregates exposure to the settlement across
desks, limits trading in the underlying by whoever benefits from its effect
on a derivative, and documents the reasons \emph{before} the trade.
\end{solution}

\begin{solution}{pb:m1:options-markets-europe-asia:1}
\textbf{1.} $90/0.85 = 106$ to 1.
\textbf{2.} M1: $90 \times 10^9 \times 6\,600 = \$594$ trillion. M2: $0.85 \times
10^9 \times 560\,000 = \$476$ trillion.
\textbf{3.} M1: $90 \times 10^9 \times 10.5 = \$945$ billion. M2: $0.85 \times 10^9
\times 1\,600 = \$1\,360$ billion.
\textbf{4.} 1.44.
\textbf{5.} 15.9 and 28.6 basis points. M1's options are cheaper relative to
their underlying: shorter-dated and further out of the money, lottery
tickets on the day's move.
\textbf{6.} M1: $945 \times 1\% = \$9.45$ billion. M2: $1\,360 \times 0.375\% =
\$5.10$ billion.
\textbf{7.} Fees $0.05\% \times 945 = \$0.47$ billion; tax on the half of
premium that makers sell, $0.1\% \times 472.5 = \$0.47$ billion: \$0.94 billion.
\textbf{8.} $0.85 \times 10^9 \times 0.50 = \$0.42$ billion.
\textbf{9.} \$8.51 billion and \$4.67 billion.
\textbf{10.} A small one at first. It depends on access (membership,
co-location, market-maker status and its privileges), on speed and pricing
quality relative to incumbents, and on which part of the flow it sees: the
pool's average spread is earned on benign flow that incumbents with
entitlements see first.
\textbf{11.} $24 \times 75/(90 \times 25) = 80\%$.
\textbf{12.} $24 \times 10^9 \times 75 \times 35 \times 0.012 = \$756$ billion.
\textbf{13.} $756 \times (1\% - 0.1\%) = \$6.8$ billion, a fifth less; and
the participants who left were the least informed, so spreads earned on the
remainder are likely to be lower too.
\textbf{14.} A business built on one market's retail flow has a policy risk
that no model hedges: the regulator can change lot, expiries, margins or
access in a quarter. Size fixed costs to the market that would survive such
a change, and hold a second market.
\textbf{15.} Its study found 93\% of individual traders losing money, in
aggregate more than 1.8 trillion rupees in three years, much of it in
expiry-day options. Larger contracts, fewer weekly expiries, upfront
premium and extra expiry-day margin all raise the price of the cheapest
bet.
\textbf{16.} Transaction costs, brokers and exchanges first; then the
counterparties with better pricing and execution: proprietary firms and
foreign investors, largely algorithmic.
\textbf{17.} No: quoting two-sided prices to willing buyers is the service.
The line is crossed when a participant trades the underlying in order to
move the settlement on which its options depend; the proposition shows that
the temptation is structural wherever the options market dwarfs its
underlying, which is why the surveillance is too.
\textbf{18.} A legal route (registration category, local entity or
partner), a clearing member, tax structuring, \omterm{def:m1:options-markets-europe-asia:weekly}{position limits} that fit the
strategy, local compliance staff, and an understanding of how the regulator
views foreign proprietary trading.
\textbf{19.} \textbf{1.44: the market with 1\% of the contracts has 44\% more
\omterm{def:m1:options-markets-europe-asia:turnover}{premium turnover}.}
\textbf{20.} The number of times a matching engine printed a trade, which is
the exchange's revenue driver and nobody else's.
\end{solution}

\begin{solution}{iq:m1:options-markets-europe-asia:1}
By what measure: contracts, notional or premium? What is the \omterm{def:m1:options-markets-europe-asia:turnover}{lot size} and
how has it changed? How much is options against futures, index against
single-stock, how much expires the same day? Who trades: the individual
share, and the share of proprietary firms? And what would the ranking be in
\omterm{def:m1:options-markets-europe-asia:turnover}{premium turnover}, which is the money.
\iqlookfor{the three measures without prompting.}
\end{solution}

\begin{solution}{iq:m1:options-markets-europe-asia:2}
One exchange and one clearing house per product, not many exchanges on one
clearing house: no fungibility across venues, no \omterm{def:m1:us-equity-market-structure:nbbo}{national best bid and offer},
no \omterm{def:m1:options-market-structure:auction}{price-improvement auctions}. Index options are European-style and cash
settled, hedged against the local index future. Much institutional size
trades as negotiated blocks. Retail demand for leverage goes largely to
bank-issued warrants and certificates, not to listed options.
\iqlookfor{the clearing structure as the root difference.}
\end{solution}

\begin{solution}{iq:m1:options-markets-europe-asia:3}
Contracts fall by at least the factor of three mechanically. Notional and
premium fall less, by however many participants are priced out: the smallest
accounts, whose minimum ticket has tripled. Remaining flow is somewhat
better capitalised and better informed; spreads in premium terms may
tighten. \omterm{def:m1:what-a-trading-firm-does:market-maker}{Market makers} lose the most benign flow; exchanges lose fee
revenue per unit of notional if fees are per contract. Some activity
migrates to other leveraged products.
\iqlookfor{mechanical versus behavioural effects.}
\end{solution}

\begin{solution}{iq:m1:options-markets-europe-asia:4}
Their payoff depends on a number computed from a few prices in a short
window, in a market that may be far smaller than the derivative positions
referencing it, so moving the underlying slightly can be worth more than it
costs. Defences: settlement on an average over a window or on an auction
with imbalance information and price collars; \omterm{def:m1:options-markets-europe-asia:weekly}{position limits} monitored
intraday on expiry day; higher margins on expiry day; surveillance that
joins participants' cash and derivatives activity; and spreading expiries
away from single thin underlyings.
\iqlookfor{the size mismatch and at least three defences.}
\end{solution}

\begin{solution}{iq:m1:options-markets-europe-asia:5}
A firm-wide view, in real time, of exposure to the settlement value across
all desks and entities. A rule that whoever benefits from the index moving
may not trade the constituents aggressively in the settlement window, or
only through a pre-approved schedule. Pre-trade documentation of the reason
for large cash trades on expiry day. Limits on the firm's share of volume in
the constituents during the window. Independent review of expiry-day \omterm{def:m1:pnl-and-positions:attribution}{P\&L
attribution}. And escalation when an exchange or regulator asks questions.
\iqlookfor{aggregation across desks, and documentation before the trade.}
\end{solution}

\begin{solution}{iq:m1:options-markets-europe-asia:6}
Size in money: \omterm{def:m1:options-markets-europe-asia:turnover}{premium turnover}, spreads, hence the revenue pool; then fees
and taxes. Flow: who trades, how informed, how concentrated in time
(expiry days). Access: legal form, membership, co-location, market-maker
programmes and their obligations, \omterm{def:m1:options-markets-europe-asia:weekly}{position limits}. Competition: incumbents'
share and edge. Policy: the regulator's stated concerns, recent and pending
changes to lots, expiries and margins, attitude to foreign proprietary
firms; scenario with the pool cut by half. Operations: clearing, collateral,
currency, tax, staff. Exit: fixed costs recoverable if the rules change.
\iqlookfor{policy risk as a first-class item, with a scenario.}
\end{solution}
